\section{Neural Compute、Memory 与 Personalization}

完成行动演示后，课堂回到架构层，解释模型在系统中的位置以及 agent 怎样跨步骤、跨会话保存信息。这里最容易混淆的是 context、persistent memory 与 user profile：三者都向模型提供信息，但生命周期、可信度和隐私义务不同。

\subsection{Model as Neural Compute Unit}

Neural Compute Unit slide 把 Transformer 画成输入输出之间的计算核心，并用 CPU analogy 强调它可以被更大系统调用。类比有助于理解 modularity，但 CPU 执行确定指令，语言模型生成概率分布；因此 agent 还需要外部规则和 verifier 约束不确定输出。

\lecturefigure{slide-51-neural-compute-unit.jpg}{AI model 作为 neural compute unit 的课堂类比}{Stanford 课堂视频 00:35:15}

\begin{knowledgebox}{读图：输入输出箭头之外还缺什么}
图中 model 接收 instructions 与 context，返回 response。真实 agent 还需 scheduler 决定何时调用、memory layer 选择哪些上下文、tool layer 执行动作、policy layer 检查权限、verifier 判断结果。模型是 compute unit，不是完整 operating system。
\end{knowledgebox}

\teachervoice{讲者反复使用 CPU/chip 类比：一个强模型只是强大的计算单元，围绕它的 abstraction、memory、interface 和 internet access 才让系统可用。这个观点贯穿后面的 LLM OS 图。}

\subsection{Context 与 Long-term Memory}

Long-term Memory slide 把 memory 分成工作记忆与持久存储。Context window 类似当前任务的短期工作区，受 token budget 限制；persistent store 保存跨会话事实、过去动作或文档索引，需要检索、更新、删除和 provenance 管理。

\lecturefigure{slide-52-long-term-memory.jpg}{Agent 的工作记忆与长期记忆问题}{Stanford 课堂视频 00:35:33}

\begin{importantbox}{Retrieval memory 的基本流程}
对记忆条目 $m_i$ 计算 embedding $e_i$，查询 $q$ 的向量为 $e_q$，可按
\[
\operatorname{score}(m_i\mid q)
=\alpha\,\cos(e_q,e_i)
+\beta\,\operatorname{recency}(m_i)
+\gamma\,\operatorname{trust}(m_i)
\]
排序。相似度回答“相关吗”，recency 处理过期信息，trust 记录来源可靠性。检索结果仍需进入 context 后由模型使用，向量库本身不会保证事实正确。
\end{importantbox}

\begin{warningbox}{Context 不是持久记忆，模型权重也不是用户数据库}
上下文窗口会在请求结束后消失或被截断，持久记忆则需要显式存储。把用户秘密写入提示、日志或无删除机制的向量库，会制造数据生命周期风险。模型权重中的训练知识也不能承担可精确更正的个人记录功能。
\end{warningbox}

\teachervoice{课堂用直观说法区分：context 是模型当下能看到的 short-term memory，数据库或其他 store 才是 long-term memory。讲者还指出长任务会迅速遇到 context limit，因此 memory management 是系统问题。}

\subsection{Personalization：显式偏好、隐式行为与反馈}

Personalization slide 提出 user agent 可学习 preferences，并用过往 actions 改进后续行为。显式偏好来自用户直接设置，例如首选航空公司；隐式偏好来自点击、选择和修改记录。后者噪声更大，也更容易越过用户预期。本节从信号来源出发，说明偏好如何进入决策，以及为什么必须允许用户查看、纠正和删除。

\lecturefigure{slide-53-personalization.jpg}{从历史行动与反馈学习个性化偏好}{Stanford 课堂视频 00:37:39}

\begin{knowledgebox}{术语消化：三类个性化信号}
\begin{tabularx}{\textwidth}{L{0.18\textwidth}X X}
\toprule
信号 & 示例 & 使用原则 \\
\midrule
Explicit & “只选直飞”“不要自动购买” & 作为高优先级约束，可查看和编辑 \\
Implicit & 多次选择某航空公司或鞋款 & 只作 soft preference，避免从一次行为过度推断 \\
Feedback & 接受、纠正、撤销或投诉 & 记录具体任务语境，区分偏好变化与系统错误 \\
\bottomrule
\end{tabularx}
\end{knowledgebox}

若推荐候选 $x$ 的基础任务效用为 $U_{task}(x)$，个性化偏好得分为 $U_{pref}(x)$，风险为 $C_{risk}(x)$，可以写成
\[
U(x)=U_{task}(x)+\lambda U_{pref}(x)-\mu C_{risk}(x).
\]
$\lambda$ 不应大到覆盖 hard constraints，$\mu$ 应随金钱、隐私和不可逆性提高。Personalization 的目标是减少重复说明，不是替用户重写目标。

\teachervoice{讲者举买鞋与机票偏好的例子，并设想从 human feedback 持续学习。讲义补上必要边界：用户点击可能只是当时妥协，不能永久固化为偏好；agent 应允许查看、纠正和遗忘。}

\subsection{Challenges：持续学习伴随隐私与控制问题}

Challenges slide 汇总 collecting user preferences、active-learning preferences、learning from interactions、on-the-fly adaptation 与 privacy。它说明 personalization 不是单纯提高推荐准确率，而是数据采集、在线更新和用户控制的综合系统问题。

\lecturefigure{slide-54-agent-challenges.jpg}{Long-term personalization 的核心挑战}{Stanford 课堂视频 00:39:54}

\begin{warningbox}{个性化系统的四个失败面}
\begin{enumerate}
  \item \textbf{错误记忆}：一次误操作被长期保存；
  \item \textbf{越权推断}：从行为推断敏感属性或跨场景复用；
  \item \textbf{陈旧偏好}：用户变化后系统仍按旧 profile 行动；
  \item \textbf{不可逆自动化}：模型把 soft preference 直接变成购买或外发动作。
\end{enumerate}
缓解措施包括 provenance、过期策略、用户可见编辑、purpose limitation 与高风险 confirmation。
\end{warningbox}

\teachervoice{讲者在展望 personalization 时主动提出 privacy concerns，并把 autonomous irreversible actions 视为危险点。课堂并没有把“知道用户更多”当作无条件目标，而是把它和权限、控制一起讨论。}

\subsection{本章小结}

Neural compute unit 类比明确了模型只是计算核心；context 提供短期工作区，persistent store 提供跨会话状态，personalization 把显式偏好、隐式行为和反馈转成可控约束。越长期的记忆越需要 provenance、更新、删除、权限与确认机制。

